\documentclass[sigconf]{acmart}
\usepackage{amsmath}
\usepackage{booktabs}
\usepackage{hyperref}
\title{Semantic-Aware Joint Source-Channel Coding for Ultra-Low-Latency Mobile Edge Networks}
\author{Fernando May Fuentes}
\affiliation{\institution{Instituto Politecnico Nacional, UPIITA}\city{Mexico City}\country{Mexico}}
\affiliation{\institution{Beihang University, School of Information and Communication Engineering}\city{Beijing}\country{China}}
\email{fmayf1500@alumno.ipn.mx}
\thanks{ORCID: \href{https://orcid.org/0009-0002-3953-5224}{0009-0002-3953-5224}.}
\begin{document}
\begin{abstract}
Mobile edge networks require a trade-off between communication cost and task accuracy. This paper presents a transparent semantic joint source-channel coding prototype that compresses 64-dimensional synthetic task features into an 8-dimensional latent representation and transmits it over a Rayleigh fading channel. The prototype is evaluated against a matched raw-feature baseline at 0, 2, 5, and 10 dB SNR. At 2 dB, the semantic representation uses 12.5\% of the raw feature dimension and achieves 90.83\% nearest-centroid task accuracy, while the raw baseline achieves 100\%. At 0 dB, semantic accuracy falls to 83.89\%. These results show a measurable bandwidth--accuracy trade-off, not a universal advantage. The artifact is a reproducible baseline for a subsequent trained Deep JSCC evaluation with public image/task data.
\end{abstract}
\keywords{semantic communication, joint source-channel coding, Rayleigh fading, mobile edge computing, 6G}
\maketitle

\section{Introduction}
Conventional communication stacks transport bits without knowing whether the receiver needs exact reconstruction or only a task-relevant representation. Semantic communication addresses this mismatch, but a useful evaluation must compare semantic compression against a matched raw baseline under the same channel. This work provides such a prototype for mobile edge networking.

\section{System Model}
An input feature vector $x\in\mathbb{R}^{64}$ is standardized and projected into a latent representation $z\in\mathbb{R}^{8}$. The Rayleigh channel is:
\begin{equation}
\hat z=h\odot z+n,\qquad n\sim\mathcal{N}(0,\sigma_n^2 I),
\end{equation}
where $h$ is sampled from a Rayleigh distribution and noise power is selected for the target SNR. The decoder applies the inverse projection and a nearest-centroid task classifier.

\section{Methods}
The prototype uses PCA as an explicit semantic compression baseline. PCA is not a neural Deep JSCC model; it provides an interpretable latent bottleneck for testing the network/channel protocol before expensive learned training. The raw baseline transmits all 64 standardized features through the same channel. The semantic baseline transmits 8 features, yielding a dimension ratio of $8/64=0.125$.

\section{Experimental Protocol}
We generate 1,200 labeled synthetic vectors from four class centers with Gaussian variation. The data are split using stratification and seed 20260915. The train/test split contains 840 and 360 samples. Experiments use SNR values 0, 2, 5, and 10 dB. The metric is nearest-centroid classification accuracy after channel corruption. This is a methodological proxy, not a real image or clinical dataset.

\section{Results}
\begin{table}[h]
\centering
\caption{Seeded semantic JSCC prototype results.}
\begin{tabular}{lccc}
\toprule
SNR & Semantic accuracy & Raw accuracy & Dimension ratio \\
\midrule
0 dB & 83.89\% & 100.00\% & 12.5\% \\
2 dB & 90.83\% & 100.00\% & 12.5\% \\
5 dB & 96.39\% & 100.00\% & 12.5\% \\
10 dB & 97.78\% & 100.00\% & 12.5\% \\
\bottomrule
\end{tabular}
\end{table}

The prototype demonstrates compression with graceful but measurable accuracy loss at low SNR. It does not reproduce the earlier unverified claim of 72\% bandwidth reduction and 95.4\% task accuracy; those claims require a trained model and public task dataset.

\section{Limitations and Future Work}
The dataset is synthetic, PCA is not Deep JSCC, the classifier is simple, and latency is not measured on a mobile-edge hardware platform. Future work will use public image/task data, a learned encoder/decoder, explicit power normalization, Rayleigh and blockage traces, repeated seeds, and latency/energy measurements.

\section{Conclusion}
This work establishes a reproducible semantic JSCC baseline for INCC 2026 and makes the compression--accuracy trade-off measurable. Its principal value is methodological: it prevents a semantic communication claim from being made without a matched channel baseline. A trained model and public-data experiment are required for the final submission.

\begin{thebibliography}{99}
\bibitem{strinati} E. C. Strinati and S. Barbarossa, ``6G networks: Beyond Shannon towards semantic and goal-oriented communications,'' Computer Networks, 2021.
\bibitem{xie} H. Xie et al., ``Deep learning enabled semantic communication systems,'' IEEE Transactions on Signal Processing, 2021.
\bibitem{bourtsoulatze} E. Bourtsoulatze et al., ``Deep joint source-channel coding for wireless image transmission,'' IEEE TCCN, 2019.
\bibitem{goldsmith} A. Goldsmith, Wireless Communications, Cambridge University Press, 2005.
\end{thebibliography}
\end{document}
